% SHARD-ID: CA-47-QUBIT-PAULI-WORD-MOMENT-BASIS
% SHARD-TITLE: Pauli Word Moment Basis for Infinite Qubit Chains
% SHARD-SUMMARY: Implements finite positioned Pauli words and translation-canonical moment variables.
% SHARD-SUMMARY: Exact Pauli multiplication gives the phases used in moment matrices and residual constraints.
% SHARD-SUMMARY: Rows and columns are actual local words; only the moment variables are quotiented by translation.
% SHARD-KEYWORDS: Pauli words, moment basis, translation invariance, qubit chain, exact algebra

\section{Pauli Word Moment Basis for Infinite Qubit Chains}
\label{sec:qubit-pauli-word-moment-basis}

\subsection{Status, Sources, and Dependencies}

\textbf{Status: checked.}  The word algebra is implemented without dense
matrices and checked against exact Pauli identities.

Dependencies: CA-38, CA-43, CA-46.

% Source: references/text/CFTFromLatticeFermions.txt:80--100 -- local finite
%   support Hamiltonian terms and observable algebra.
% Source: references/text/CFTFromLatticeFermions.txt:1033--1038 -- Pauli
%   matrices in the qubit setting.
% Checked: src/QubitPauliWords.jl; test/runtests.jl testset
%   "qubit Pauli moment SDP hierarchy".

\subsection{Positioned Words}

A word is stored as sorted non-identity labels on integer sites.  Identities
are not stored, but gaps are preserved:
\[
  X_0Z_2\quad\leftrightarrow\quad
  \texttt{sites=(0,2), labels=(1,3)}.
\]
The empty word is the identity.  Multiplication is sitewise with
\[
  XY=iZ,\quad YZ=iX,\quad ZX=iY,
\]
and the reversed products carrying \(-i\).  The implementation returns
\[
  P_uP_v=\phi(u,v)P_{u\cdot v},
  \qquad \phi(u,v)\in\{1,-1,i,-i\}.
\]
The tests assert \(XY=iZ\), \(YX=-iZ\), \(X^2=I\), and commutation of disjoint
words.

\subsection{Translation Quotient}

The SDP row/column words are not quotiented by translation.  They are actual
local probes in a finite window.  Translation invariance enters only when a
moment variable is named:
\[
  y_s=\omega(P_s)=\omega(P_{\tau^n s}).
\]
The canonical representative shifts the first non-identity site to \(0\).
Thus \(X_{-3}Z_{-1}\) and \(X_0Z_2\) share the same moment variable.

Because each phase-free Pauli word is self-adjoint, each canonical moment
variable is real.  Complex numbers enter only as multiplication phases.
